\documentclass[11pt]{article}
\usepackage[letterpaper,margin=1in]{geometry}
\usepackage{amsmath,amssymb,amsthm,booktabs,microtype,xurl}
\usepackage[hidelinks]{hyperref}
\newtheorem{theorem}{Theorem}
\newtheorem{lemma}[theorem]{Lemma}
\newcommand{\C}{\mathbb C}
\newcommand{\Q}{\mathbb Q}
\newcommand{\N}{\mathbb N}
\newcommand{\SIC}{\operatorname{SIC}}

\title{An explicit counterexample to the\\Special Image Conjecture in dimension 21}
\author{Alec Kriebel\\[2pt]\large with heavy assistance from ChatGPT 5.6 Sol}
\date{First public branch draft 22 July 2026, 02:59:33 UTC\\
Research draft; not peer reviewed}

\begin{document}
\maketitle

\begin{abstract}
Let
\[
 E:\C[\xi_1,\ldots,\xi_n,Z_1,\ldots,Z_n]\longrightarrow\C[Z],
 \qquad E(\xi^\alpha p(Z))=\partial_Z^\alpha p(Z).
\]
The Special Image Conjecture $\SIC(n)$ says that $\ker E$ is a
Mathieu--Zhao subspace. We give an exact rational polynomial
$A\in\Q[\xi,Z]$ for $n=21$, with 72 monomials and total degree four,
and take $b=Z_1$. We prove
\[
 A^m\in\ker E\quad(m\geq1),\qquad
 bA^m\notin\ker E\quad\text{for infinitely many }m.
\]
Thus $\SIC(21)$ is false. The witness is obtained from a certified
13-variable stable model of the newly announced Keller counterexample by
adjoining eight variables. The bridge is a scalar-parameter form of
Abhyankar--Gurjar inversion; unlike Zhao's homogeneous specialization, it
permits the linear block in this construction. Exact SymPy and
dependency-free rational checkers accompany a 72-term sparse certificate.
\end{abstract}

\noindent\textbf{Verification disclaimer.}
The human author is a complete amateur exploring the limits of AI-assisted
mathematics and cannot independently verify these claims. This is a research
artifact for expert checking, not an established result. Priority is
provisional and limited by the documented audit.

\section{The Special Image Conjecture}

Put $R_n=\C[\xi_1,\ldots,\xi_n,Z_1,\ldots,Z_n]$ and define
$E_n:R_n\to\C[Z]$ by
\[
 E_n(\xi^\alpha p(Z))=\partial_Z^\alpha p(Z).
\]
If $\Theta_j=\xi_j-\partial_{Z_j}$, Zhao's theorem gives
\begin{equation}
 \ker E_n=\sum_{j=1}^n\Theta_jR_n. \label{eq:kernel}
\end{equation}
A linear subspace $M$ of a commutative algebra is a Mathieu--Zhao subspace if
\[
 f^m\in M\ (m\geq1)
 \quad\Longrightarrow\quad
 qf^m\in M\ (m\gg0)\quad\text{for every }q.
\]
The assertion that $\ker E_n$ is such a subspace is $\SIC(n)$.

\section{A scalar-parameter inversion lemma}

The following form is useful because it does not require $g$ to be
homogeneous or to have order at least two.

\begin{lemma}[pencil-Keller collision lemma]\label{lem:pencil}
Let $k$ have characteristic zero, let $g\in k[Z]^n$, and suppose
\begin{equation}
 \det(I+tJg)=1\quad\text{in }k[t,Z]. \label{eq:pencil}
\end{equation}
Set
\[
 A(\xi,Z)=-\sum_{j=1}^n\xi_jg_j(Z).
\]
Then $E_n(A^m)=0$ for every $m\geq1$. If $T=I+g$ has two points $p,q$
with $T(p)=T(q)$ and $p_i\ne q_i$, then
\[
 E_n(Z_iA^m)\ne0
\]
for infinitely many $m$.
\end{lemma}

\begin{proof}
Work in the $t$-adically complete ring $k[Z][[t]]$. The map
\[
 T_t(Z)=Z+tg(Z)=Z-H_t(Z),\qquad H_t=-tg,
\]
has a unique formal inverse $Q_t$, because it is congruent to the identity
modulo $t$. The Abhyankar--Gurjar inversion identity is valid in this
filtration: for $p\in k[Z]$,
\begin{equation}
 p(Q_t(Z))=
 \sum_{\alpha\in\N^n}\frac1{\alpha!}\partial_Z^\alpha
 \bigl(H_t(Z)^\alpha\det JT_t(Z)\,p(Z)\bigr). \label{eq:AG}
\end{equation}
Indeed, apply the usual formal identity modulo $t^{N+1}$ for every $N$;
$H_t^\alpha$ has $t$-adic order $|\alpha|$, so every coefficient is a finite
sum. By~\eqref{eq:pencil}, $\det JT_t=1$. The multinomial theorem turns
\eqref{eq:AG} into
\begin{equation}
 p(Q_t(Z))=\sum_{m\geq0}\frac{t^m}{m!}E_n\bigl(p(Z)A^m\bigr). \label{eq:coeff}
\end{equation}
Taking $p=1$ makes the left side equal to one, proving
$E_n(A^m)=0$ for every positive $m$.

Now take $p=Z_i$. If $E_n(Z_iA^m)=0$ for all sufficiently large $m$,
\eqref{eq:coeff} says that the $i$th coordinate of $Q_t$ is a polynomial in
$t,Z$. The formal identity $Q_{t,i}(T_t(Z))=Z_i$ is then a polynomial
identity and may be specialized at $t=1$. Hence
$Q_{1,i}(T(p))=p_i$ and $Q_{1,i}(T(q))=q_i$, contradicting
$T(p)=T(q)$ and $p_i\ne q_i$.
\end{proof}

The $t$-adic formulation matters below: $g$ has a linear $BU$ block. Zhao's
homogeneous theorem should therefore not be quoted verbatim for this witness;
Lemma~\ref{lem:pencil} supplies the needed extension.

\section{The explicit 21-variable map}

Use the coordinate order
\[
 Z=(x,y,z,a_1,b_1,a_2,b_2,a_3,b_3,a_4,b_4,a_5,b_6,U_1,\ldots,U_8).
\]
Define $g=(g_1,\ldots,g_{21})$ as follows.

\begingroup\small
\begin{align}
g_1={}&U_1+\tfrac12a_1b_1+\tfrac32a_1y,&
g_2={}&U_2-a_2b_2-2a_2z+a_3b_3+3xz,\nonumber\\
g_3={}&U_3+b_1a_5-a_3b_6-a_4b_4-7a_4y+3a_5y+4y^2,&
g_4={}&x^2,\nonumber\\
g_5={}&-3U_2+3a_2b_2+6a_2z-3a_3b_3-8xz,&
g_6={}&U_4,\nonumber\\
g_7={}&U_5-2b_1a_5+2a_3b_6+2a_4b_4+14a_4y-6a_5y-5y^2,&
g_8={}&xy,\nonumber\\
g_9={}&a_2z+3b_2x,&g_{10}={}&U_6,\nonumber\\
g_{11}={}&U_7-b_1a_3+7a_2b_2+14a_2z-7a_3b_3-3a_3y-18xz,&
g_{12}={}&U_8,\nonumber\\
g_{13}={}&b_1a_4-b_4y.&& \label{eq:gfirst}
\end{align}
\begin{align}
g_{14}={}&-\tfrac12a_1xz-\tfrac12b_1x^2,\nonumber\\
g_{15}={}&-a_2a_3z+3a_2y^2-3b_2a_3x-b_3xy-12xy^2,\nonumber\\
g_{16}={}&b_1a_3a_4-a_3b_4y+3a_4xz-a_5xz+b_6xy-3xyz,\nonumber\\
g_{17}={}&-3x^2y,\nonumber\\
g_{18}={}&-2b_1a_3a_4+2a_3b_4y-6a_4xz+2a_5xz-2b_6xy+5xyz,\nonumber\\
g_{19}={}&-xy^2,\nonumber\\
g_{20}={}&b_1xy+7a_2a_3z-21a_2y^2+21b_2a_3x+a_3xz+7b_3xy+84xy^2,\nonumber\\
g_{21}={}&-a_4xy. \label{eq:gsecond}
\end{align}
\endgroup

The polynomial in the theorem is
\begin{equation}
 A(\xi,Z)=-\sum_{j=1}^{21}\xi_jg_j(Z),\qquad b=Z_1=x. \label{eq:A}
\end{equation}
It has exactly 72 monomials and total degree four. Equations
\eqref{eq:gfirst}--\eqref{eq:A} are a human-readable certificate; the
accompanying JSON records the same object as exact sparse rational data.

\section{Determinant certificate}

The structure behind the display is an exactly certified 13-variable stable
model from Exploration 03,
\begin{equation}
 \Psi(X)=X+H_2(X)+BK(X), \label{eq:Psi}
\end{equation}
where $H_2$ is quadratic, $K$ has eight cubic components, and
$B\in\operatorname{Mat}_{13\times8}(\Q)$. Its Jacobian determinant is one.
With $Z=(X,U)$, the displayed vector is
\begin{equation}
 g(X,U)=\bigl(H_2(X)+BU,-K(X)\bigr). \label{eq:blockg}
\end{equation}
Therefore
\[
 I+sJg=\begin{pmatrix}I_{13}+sJH_2&sB\\-sJK&I_8\end{pmatrix}.
\]
Taking the Schur complement of $I_8$, and using the homogeneities of $H_2$
and $K$, gives the polynomial identity
\begin{align}
\det(I+sJg)
 &=\det\bigl(I_{13}+sJH_2+s^2BJK\bigr)\nonumber\\
 &=\det J\Psi(sX)=1. \label{eq:detcert}
\end{align}
The primary checker verifies the matrix equality in~\eqref{eq:detcert}
entry by entry. The dependency-free checker separately tests the determinant
pencil at 66 exact rational specializations; a Node.js/BigInt implementation
supplies a third runtime and checks 18 further exact specializations.

\section{Collision certificate and theorem}

The following three rational points use the coordinate order above.

\begingroup\small
\begin{align*}
p_0={}&(0,0,-\tfrac14,0,0,0,\tfrac12,0,0,0,0,0,0;
0,0,0,0,0,0,0,0),\\
p_+={}&(1,-\tfrac32,\tfrac{13}2,-1,-2,\tfrac92,-10,\tfrac32,\tfrac34,
-\tfrac94,-6,-\tfrac{27}8,\tfrac92;\\
&\hspace{8mm}-\tfrac{17}4,-\tfrac{45}8,\tfrac{99}{16},-\tfrac92,
-\tfrac{177}8,\tfrac94,\tfrac{213}8,\tfrac{27}8),\\
p_-={}&(-1,\tfrac32,\tfrac{13}2,-1,2,-\tfrac92,-10,\tfrac32,-\tfrac34,
\tfrac94,6,\tfrac{27}8,\tfrac92;\\
&\hspace{8mm}\tfrac{17}4,\tfrac{45}8,\tfrac{99}{16},\tfrac92,
-\tfrac{177}8,-\tfrac94,-\tfrac{213}8,-\tfrac{27}8).
\end{align*}
\endgroup
Direct substitution gives
\begin{equation}
 (I+g)(p_0)=(I+g)(p_+)=(I+g)(p_-)
 =(0,0,-\tfrac14,0,0,0,\tfrac12,0,\ldots,0). \label{eq:collision}
\end{equation}
Their first coordinates are $0,1,-1$.

\begin{theorem}\label{thm:main}
For $A$ and $b$ in~\eqref{eq:A},
\[
 A^m\in\ker E_{21}\quad\text{for every }m\geq1,
\]
but
\[
 bA^m\notin\ker E_{21}
\]
for infinitely many $m$. Consequently $\SIC(21)$ is false.
\end{theorem}

\begin{proof}
Equation~\eqref{eq:detcert} and the collision $p_0,p_+$ in
\eqref{eq:collision}, separated by their first coordinate, satisfy
Lemma~\ref{lem:pencil} with $n=21$ and $i=1$. The two assertions follow.
Equation~\eqref{eq:kernel} identifies $\ker E_{21}$ with the image appearing
in the usual statement of the Special Image Conjecture.
\end{proof}

As low-order guards, the symbolic checker also obtains
\begin{equation}
 E_{21}(A)=E_{21}(A^2)=0,\quad E_{21}(bA)\ne0,\quad
 E_{21}(bA^2)=3x^2y\ne0. \label{eq:guards}
\end{equation}

\section{Verification and scope}

The repository contains three exact implementations. The SymPy
checker reconstructs the stable model, verifies the Schur-complement matrix
identity, the 72-term count, the collision, and~\eqref{eq:guards}. A
dependency-free checker parses only the exported JSON, rebuilds $A=-\xi\cdot
g$, evaluates the collision with rational arithmetic, and checks the
determinant pencil at 66 exact specializations. A separate Node.js/BigInt
implementation rechecks the collision and 18 determinant specializations
without Python or a computer algebra system.

Zhang's immediate-consequences note correctly established only that the Image
Conjecture fails in some finite dimension. There is nevertheless an immediate
explicit predecessor: Exploration 03's 22-variable cubic homogeneous map,
combined with Zhao's theorem, gives an $\SIC(22)$ witness. Thompson's
24-variable and Harrison's 79-variable cubic models similarly imply larger
witnesses, whether or not those corollaries were written out as pairs.

The contribution here is narrower: removing the homogenizing variable lowers
the explicit dimension from 22 to 21. The resulting linear $BU$ block requires
Lemma~\ref{lem:pencil}. At the audit cutoff, no earlier public $\SIC(n)$
witness with $n\leq21$, or the nonhomogeneous lemma in the form used here, was
located.

This is evidence, not proof, of priority. Public work is changing hourly and
private or unindexed work may exist. We claim neither minimality of dimension
nor minimality of support.

\begin{thebibliography}{9}
\bibitem{Zhao2010}
W.~Zhao,
\emph{Images of commuting differential operators of order one with constant
leading coefficients}, J. Algebra 324 (2010), 231--247;
\url{https://arxiv.org/abs/0902.0210}.

\bibitem{DEZ2017}
H.~Derksen, A.~van den Essen, and W.~Zhao,
\emph{The Gaussian Moments Conjecture and the Jacobian Conjecture},
Israel J. Math. 219 (2017), 917--928;
\url{https://arxiv.org/abs/1506.05192}.

\bibitem{Zhang2026}
Z.~Zhang,
\emph{Direct Consequences of the Three-Dimensional Counterexample to the
Jacobian Conjecture}, 20 July 2026;
\url{https://zzhang-iu.github.io/papers/direct-consequences-jacobian/}.

\bibitem{KriebelD3}
A.~Kriebel, with ChatGPT 5.6 Sol,
\emph{Small explicit symmetric Keller and vanishing counterexamples},
Exploration 03, 21 July 2026.
\end{thebibliography}

\end{document}
